\textbf{Phần 1. Trắc nghiệm lựa chọn}
\setcounterpageref{ex}{0}

\begin{ex} %[3H1V17]
Một hình tròn có bán kính bằng $14\text{ cm}$. Đường kính của hình tròn đó là
\choice
{$7\text{ cm}$}
{$21\text{ cm}$}
{\True $28\text{ cm}$}
{$35\text{ cm}$}
\loigiai{
Đường kính gấp 2 lần bán kính: $14 \times 2 = 28\text{ (cm)}$.
\textbf{Đáp án của câu hỏi là C}
}
\end{ex}

\begin{ex} %[3H1V17]
Một hình tròn có đường kính dài bằng chu vi của hình vuông cạnh $12\text{ cm}$. Bán kính của hình tròn đó bằng
\choice
{\True $24\text{ cm}$}
{$48\text{ cm}$}
{$12\text{ cm}$}
{$36\text{ cm}$}
\loigiai{
Chu vi hình vuông là: $12 \times 4 = 48\text{ (cm)}$.\
Đường kính hình tròn là $48\text{ cm}$ nên bán kính là: $48 : 2 = 24\text{ (cm)}$.
\textbf{Đáp án của câu hỏi là A}
}
\end{ex}

\begin{ex} %[3H1V17]
Cho hình tròn tâm $O$ có bán kính $R = 15\text{ cm}$. Điểm $A$ nằm trong hình tròn, điểm $B$ nằm trên đường tròn và điểm $C$ nằm ngoài hình tròn. Khẳng định nào sau đây là đúng?
\choice
{$OA = 15\text{ cm}$ và $OB > 15\text{ cm}$}
{$OA > 15\text{ cm}$ và $OC < 15\text{ cm}$}
{$OB < 15\text{ cm}$ và $OC = 15\text{ cm}$}
{\True $OA < 15\text{ cm}$ và $OC > 15\text{ cm}$}
\loigiai{
Điểm nằm trong hình tròn có khoảng cách tới tâm nhỏ hơn bán kính: $OA < 15\text{ cm}$. Điểm nằm ngoài hình tròn có khoảng cách tới tâm lớn hơn bán kính: $OC > 15\text{ cm}$.
\textbf{Đáp án của câu hỏi là D}
}
\end{ex}

\begin{ex} %[3H1V17]
Hai hình tròn có chung tâm $O$. Hình tròn thứ nhất có bán kính $18\text{ cm}$, hình tròn thứ hai có bán kính gấp đôi bán kính hình tròn thứ nhất. Đường kính của hình tròn thứ hai là
\choice
{$36\text{ cm}$}
{\True $72\text{ cm}$}
{$54\text{ cm}$}
{$90\text{ cm}$}
\loigiai{
Bán kính hình tròn thứ hai là: $18 \times 2 = 36\text{ (cm)}$.\
Đường kính hình tròn thứ hai là: $36 \times 2 = 72\text{ (cm)}$.
\textbf{Đáp án của câu hỏi là B}
}
\end{ex}

\begin{ex} %[3H1V17]
Có ba hình tròn bằng nhau đặt thẳng hàng tiếp xúc nhau sao cho tâm của chúng cùng nằm trên một đoạn thẳng. Biết khoảng cách giữa tâm hình tròn thứ nhất và tâm hình tròn thứ ba là $28\text{ cm}$. Bán kính của mỗi hình tròn là
\choice
{$14\text{ cm}$}
{$4\text{ cm}$}
{\True $7\text{ cm}$}
{$8\text{ cm}$}
\loigiai{
Khoảng cách giữa tâm thứ nhất và tâm thứ ba gồm 2 đoạn nối tâm liền kề, mỗi đoạn bằng 2 lần bán kính: khoảng cách bằng $4$ lần bán kính.\
Bán kính mỗi hình tròn là: $28 : 4 = 7\text{ (cm)}$.
\textbf{Đáp án của câu hỏi là C}
}
\end{ex}

\begin{ex} %[3H1V17]
Một compa mở hai chân tạo một khoảng cách bằng $9\text{ cm}$. Khi quay một vòng để vẽ đường tròn thì đường tròn nhận được có đường kính là
\choice
{\True $18\text{ cm}$}
{$9\text{ cm}$}
{$27\text{ cm}$}
{$36\text{ cm}$}
\loigiai{
Khoảng cách giữa hai chân compa chính là bán kính hình tròn: $R = 9\text{ cm}$.\
Đường kính của hình tròn là: $9 \times 2 = 18\text{ (cm)}$.
\textbf{Đáp án của câu hỏi là A}
}
\end{ex}

\begin{ex} %[3H1V17]
Trong một hình chữ nhật có chiều dài $20\text{ cm}$ và chiều rộng $12\text{ cm}$, người ta vẽ một hình tròn có bán kính lớn nhất có thể nằm hoàn toàn bên trong hình chữ nhật đó. Đường kính của hình tròn đó bằng
\choice
{$20\text{ cm}$}
{$6\text{ cm}$}
{$10\text{ cm}$}
{\True $12\text{ cm}$}
\loigiai{
Hình tròn lớn nhất nằm trọn trong hình chữ nhật có đường kính bằng đúng chiều rộng hình chữ nhật, tức là $12\text{ cm}$.
\textbf{Đáp án của câu hỏi là D}
}
\end{ex}

\begin{ex} %[3H1V17]
Tổng độ dài của đường kính và bán kính của một hình tròn là $45\text{ cm}$. Bán kính của hình tròn đó là
\choice
{$30\text{ cm}$}
{\True $15\text{ cm}$}
{$10\text{ cm}$}
{$20\text{ cm}$}
\loigiai{
Vì đường kính gấp đôi bán kính nên tổng độ dài đường kính và bán kính bằng 3 lần bán kính.\
Bán kính hình tròn là: $45 : 3 = 15\text{ (cm)}$.
\textbf{Đáp án của câu hỏi là B}
}
\end{ex}

\begin{ex} %[3H1V17]
Đường kính của hình tròn thứ nhất bằng bán kính của hình tròn thứ hai. Biết bán kính hình tròn thứ nhất là $8\text{ cm}$. Đường kính của hình tròn thứ hai là
\choice
{$16\text{ cm}$}
{$24\text{ cm}$}
{\True $32\text{ cm}$}
{$64\text{ cm}$}
\loigiai{
Đường kính hình tròn thứ nhất là: $8 \times 2 = 16\text{ (cm)}$.\
Bán kính hình tròn thứ hai bằng $16\text{ cm}$, suy ra đường kính của nó là: $16 \times 2 = 32\text{ (cm)}$.
\textbf{Đáp án của câu hỏi là C}
}
\end{ex}

\begin{ex} %[3H1V17]
Một hình tròn tâm $O$ có bán kính $R = 25\text{ cm}$. Đoạn thẳng $MN$ đi qua tâm $O$ và có hai đầu mút nằm trên đường tròn. Đoạn thẳng $MN$ có độ dài là
\choice
{\True $50\text{ cm}$}
{$25\text{ cm}$}
{$75\text{ cm}$}
{$100\text{ cm}$}
\loigiai{
Đoạn thẳng đi qua tâm và có hai mút nằm trên đường tròn chính là đường kính của hình tròn: $MN = 25 \times 2 = 50\text{ (cm)}$.
\textbf{Đáp án của câu hỏi là A}
}
\end{ex}

\begin{ex} %[3H1V17]
Một bánh xe đạp có bán kính $30\text{ cm}$. Khoảng cách từ trục bánh xe (tâm) đến một điểm bất kì trên vành bánh xe là
\choice
{$60\text{ cm}$}
{$15\text{ cm}$}
{$45\text{ cm}$}
{\True $30\text{ cm}$}
\loigiai{
Khoảng cách từ tâm đến một điểm bất kì trên đường tròn chính là bán kính của đường tròn, bằng $30\text{ cm}$.
\textbf{Đáp án của câu hỏi là D}
}
\end{ex}

\begin{ex} %[3H1V17]
Cho đoạn thẳng $AB = 10\text{ cm}$. Vẽ đường tròn tâm $A$ bán kính $6\text{ cm}$ và đường tròn tâm $B$ bán kính $4\text{ cm}$. Khẳng định nào sau đây là đúng?
\choice
{Hai đường tròn cắt nhau tại 2 điểm}
{\True Hai đường tròn tiếp xúc nhau tại đúng 1 điểm trên đoạn thẳng $AB$}
{Hai đường tròn không có điểm chung nào}
{Đường tròn tâm $B$ nằm hoàn toàn bên trong đường tròn tâm $A$}
\loigiai{
Tổng hai bán kính là $6 + 4 = 10\text{ cm}$, đúng bằng khoảng cách giữa hai tâm $AB = 10\text{ cm}$. Do đó hai đường tròn tiếp xúc ngoài tại một điểm trên đoạn $AB$.
\textbf{Đáp án của câu hỏi là B}
}
\end{ex}

\textbf{Phần 2. Trắc nghiệm đúng sai}
\setcounterpageref{ex}{0}

\begin{ex} %[3H1V17]
Xét tính đúng sai của các khẳng định về hình tròn:
\choiceTF
{\True Đường kính của hình tròn luôn có độ dài gấp 2 lần bán kính của hình tròn đó.}
{Tất cả các đoạn thẳng nối hai điểm trên đường tròn đều là đường kính.}
{\True Tâm của hình tròn chính là trung điểm của mọi đường kính của hình tròn đó.}
{Hai hình tròn có bán kính khác nhau có thể có cùng đường kính.}
\loigiai{
\begin{itemize}
    \item Ý a) Đúng: Đường kính $d = 2 \times R$.
    \item Ý b) Sai: Sai vì chỉ đoạn thẳng đi qua tâm mới là đường kính.
    \item Ý c) Đúng: Tâm nằm giữa và chia đôi đường kính.
    \item Ý d) Sai: Sai vì bán kính khác nhau thì đường kính cũng phải khác nhau.
\end{itemize}
\begin{itemize}
    \item a) \textbf{Đúng}.
    \item b) \textbf{Sai}.
    \item c) \textbf{Đúng}.
    \item d) \textbf{Sai}.
\end{itemize}
}
\end{ex}

\begin{ex} %[3H1V17]
Cho hình tròn tâm $O$ có đường kính $AB = 24\text{ cm}$:
\choiceTF
{\True Bán kính của hình tròn bằng $12\text{ cm}$.}
{Nếu điểm $C$ nằm ngoài hình tròn thì đoạn thẳng $OC$ có thể ngắn hơn $12\text{ cm}$.}
{\True Điểm $O$ là trung điểm của đoạn thẳng $AB$.}
{\True Mọi bán kính xuất phát từ tâm $O$ đều có độ dài bằng nhau và bằng $12\text{ cm}$.}
\loigiai{
\begin{itemize}
    \item Ý a) Đúng: $24 : 2 = 12\text{ cm}$.
    \item Ý b) Sai: Sai vì điểm ngoài hình tròn luôn có $OC > 12\text{ cm}$.
    \item Ý c) Đúng: Đoạn $AB$ là đường kính nên tâm $O$ là trung điểm.
    \item Ý d) Đúng: Tính chất của hình tròn.
\end{itemize}
\begin{itemize}
    \item a) \textbf{Đúng}.
    \item b) \textbf{Sai}.
    \item c) \textbf{Đúng}.
    \item d) \textbf{Đúng}.
\end{itemize}
}
\end{ex}

\textbf{Phần 3. Trả lời ngắn}
\setcounterpageref{ex}{0}

\begin{ex} %[3H1V17]
Một hình tròn có đường kính bằng $36\text{ cm}$. Tính bán kính của hình tròn đó theo xăng-ti-mét.
\shortans{18}
\loigiai{
\textbf{Đáp số:} $18$\\
Bán kính là: $36 : 2 = 18\text{ (cm)}$.
}
\end{ex}

\begin{ex} %[3H1V17]
Một hình tròn có bán kính $8\text{ cm}$. Hỏi đường kính của hình tròn đó dài bao nhiêu xăng-ti-mét?
\shortans{16}
\loigiai{
\textbf{Đáp số:} $16$\\
Đường kính là: $8 \times 2 = 16\text{ (cm)}$.
}
\end{ex}

\begin{ex} %[3H1V17]
Cho đoạn thẳng $MN = 40\text{ cm}$. Vẽ hình tròn nhận $MN$ làm đường kính. Tính khoảng cách từ tâm của hình tròn đó đến điểm $M$ theo xăng-ti-mét.
\shortans{20}
\loigiai{
\textbf{Đáp số:} $20$\\
Khoảng cách từ tâm đến điểm mút $M$ chính là bán kính: $40 : 2 = 20\text{ (cm)}$.
}
\end{ex}

\begin{ex} %[3H1V17]
Hai hình tròn tâm $A$ và tâm $B$ tiếp xúc ngoài với nhau. Biết bán kính hình tròn tâm $A$ là $15\text{ cm}$, bán kính hình tròn tâm $B$ là $9\text{ cm}$. Tính độ dài đoạn thẳng nối hai tâm $AB$ theo xăng-ti-mét.
\shortans{24}
\loigiai{
\textbf{Đáp số:} $24$\\
Độ dài đoạn nối hai tâm bằng tổng hai bán kính: $15 + 9 = 24\text{ (cm)}$.
}
\end{ex}

\begin{ex} %[3H1V17]
Một hình tròn có đường kính gấp $4$ lần bán kính của một hình tròn nhỏ có bán kính $5\text{ cm}$. Tính bán kính của hình tròn lớn theo xăng-ti-mét.
\shortans{10}
\loigiai{
\textbf{Đáp số:} $10$\\
Đường kính hình tròn lớn là: $5 \times 4 = 20\text{ cm}$. Bán kính hình tròn lớn là: $20 : 2 = 10\text{ cm}$.
}
\end{ex}

\begin{ex} %[3H1V17]
Người ta xếp 4 đồng xu hình tròn giống hệt nhau vừa khít vào một hàng thẳng có tổng chiều dài $24\text{ cm}$. Tính đường kính của mỗi đồng xu theo xăng-ti-mét.
\shortans{6}
\loigiai{
\textbf{Đáp số:} $6$\\
Đường kính của mỗi đồng xu là: $24 : 4 = 6\text{ (cm)}$.
}
\end{ex}
